% Zdroj: PDF priklady-leto-2023.pdf, fyzická str. 4-5. Značka: otázka studijního zaměření. Zařazeno: I4.
\section{Synchronizace}
\szzotazka{léto 2023}{14}{Synchronizace}
\noindent\textit{(Pozn.: v~rámci současné akreditace nešlo o~společnou otázku (q1--q4), ale o~otázku jiné specializace (specializační sekce, systémové zaměření, ne UI). Do společného okruhu I4 tedy formálně nepatří; tematicky sem ovšem spadá a~je výborná k~procvičení.)}\par

\noindent\textbf{Zadání.}\par
\begin{enumerate}
\item Načrtněte rozhraní semaforu a popište sémantiku jeho operací.
\item Pomocí rozhraní z prvního bodu implementujte datovou strukturu a funkce pro řešení
  problému producent-konzument, jejíž skica je uvedena níže. Není potřeba implementovat
  vlastní datovou strukturu pro udržení seznamu vyprodukovaných (ale nespotřebovaných) dat,
  stačí popsat její rozhraní. To ale samo o sobě nesmí poskytovat žádné záruky ohledně
  synchronizace.

\begin{lstlisting}[language=C]
typedef struct {
    ...
} prodcons_t;

void prodcons_init (prodcons_t *prodcons, int capacity) { ... }

void prodcons_produce (prodcons_t *prodcons, void *item) { ... }
void *prodcons_consume (prodcons_t *prodcons) { ... }
\end{lstlisting}

V~odpovědi používejte programovací jazyk C nebo C++, ale kromě rozhraní z prvního bodu
nepoužívejte žádné další nástroje pro synchronizaci.
\end{enumerate}

\medskip
\noindent\textbf{Náčrt řešení.}\par
\textbf{(1)~Rozhraní semaforu.} Semafor drží nezáporný čítač a~frontu čekajících vláken.
Operace \texttt{down} (P): je-li čítač $>0$, sníží ho; jinak vlákno \emph{pasivně} zablokuje
a~zařadí do fronty. Operace \texttt{up} (V): je-li někdo ve frontě, probudí ho; jinak čítač
zvýší. Obě operace jsou atomické.

\textbf{(2)~Producent--konzument.} Použijeme tři semafory: \texttt{free} (volná místa, init
\texttt{capacity}), \texttt{items} (připravené položky, init~0) a~binární \texttt{mtx}
(vzájemné vyloučení nad bufferem, init~1). Buffer \texttt{buf} má jen rozhraní
\texttt{buffer\_push}/\texttt{buffer\_pop} bez vlastní synchronizace.
\begin{lstlisting}[language=C]
typedef struct {
    semaphore free;    // init = capacity (volna mista)
    semaphore items;   // init = 0        (pripravene polozky)
    semaphore mtx;     // init = 1        (vzajemne vylouceni nad bufferem)
    buffer_t buf;      // push/pop bez vlastni synchronizace
} prodcons_t;

void prodcons_produce(prodcons_t *pc, void *item) {
    down(&pc->free);            // pockej na volne misto
    down(&pc->mtx);
    buffer_push(&pc->buf, item);
    up(&pc->mtx);
    up(&pc->items);             // ohlas novou polozku
}
void *prodcons_consume(prodcons_t *pc) {
    down(&pc->items);           // pockej na polozku
    down(&pc->mtx);
    void *item = buffer_pop(&pc->buf);
    up(&pc->mtx);
    up(&pc->free);              // uvolni misto
    return item;
}
\end{lstlisting}
Invarianty \texttt{free} $=\texttt{capacity}-\text{počet položek}$ a~\texttt{items} $=\text{počet
položek}$ zaručují, že se nezapisuje do plného a~nečte z~prázdného bufferu. Počítací semafor se
získává \emph{před} \texttt{mtx} (jinak by hrozil deadlock). \texttt{mtx} je potřeba i~při jediném
producentovi a~konzumentovi, protože \texttt{push} a~\texttt{pop} mohou běžet současně nad touž
nesynchronizovanou strukturou.

\smallskip
\textit{(Náčrt doplněn k~výkladu okruhu; otázka nemá v~PDF řešení, správnost ověřena rozborem invariantů a~pořadí operací.)}
